\documentclass[11pt]{article}
\usepackage[margin=1in]{geometry}
\usepackage{amsmath,amssymb,amsthm}
\usepackage{mathtools}
\usepackage{stmaryrd}

\newtheorem{theorem}{Theorem}[section]
\newtheorem{proposition}[theorem]{Proposition}
\newtheorem{lemma}[theorem]{Lemma}
\newtheorem{corollary}[theorem]{Corollary}
\theoremstyle{definition}
\newtheorem{definition}[theorem]{Definition}
\newtheorem{example}[theorem]{Example}
\theoremstyle{remark}
\newtheorem{remark}[theorem]{Remark}

\newcommand{\Z}{\mathbb{Z}}
\newcommand{\N}{\mathbb{N}}
\newcommand{\F}{\mathbb{F}}
\newcommand{\Sat}{\mathrm{Sat}}

\title{Notes On a Reverse Engineering Algorithm}
\author{}
\date{v0.1, 4 November 2026}

\begin{document}
\maketitle

\section{The machine model}

Choose a word size $w$ and let $\Z_m = \Z/m\Z$. A register file is a finite set $\mathcal{R}$ of byte cells, while memory is a total function mapping addresses to bytes. The state space is
\[
  \Sigma \;=\; \Z_{2^8}^{\mathcal{R}} \times \Z_{2^8}^{\Z_{2^{64}}}.
\]
A straight-line fragment $P$ defines a function $\llbracket P \rrbracket$ from $\Sigma$ into $\Sigma$, while an ABI defines an observation $\pi$ from $\Sigma$ into a finite set $Y$ (for instance, the low byte of the return register together with final memory). The observable function for $P$ is the composite $f_P = \pi \circ \llbracket P \rrbracket$.

All that follows is relative to this model. Faults, page permissions, timing, and concurrency are outside $\Sigma$, and none of the statements here addresses them.

\section{Boxes}

\begin{definition}
A \emph{box} is a first-order formula $A(\sigma)$ in the theory of fixed-width bit-vectors and arrays, whose free variables range over $\Sigma$. Its extension is $\Sat(A) = \{\sigma \in \Sigma : \sigma \models A\}$.
\end{definition}

A box is not a Cartesian product. For example, the formula $p = q$ over two pointer registers defines a box whose extension is the diagonal, not a rectangle. Boxes form a Boolean algebra under $\wedge$, $\vee$, $\neg$, and $\Sat$ is a homomorphism of Boolean algebras into the power set of $\Sigma$.

\begin{proposition}\label{prop:containment}
For boxes $A$ and $B$,
\[
  \Sat(A) \subseteq \Sat(B) \iff A \wedge \neg B \text{ is unsatisfiable}, \qquad
  \Sat(A) \cap \Sat(B) \neq \emptyset \iff A \wedge B \text{ is satisfiable}.
\]
\end{proposition}

\begin{proof}
$\Sat(A) \setminus \Sat(B) = \Sat(A \wedge \neg B)$, and a set is empty exactly when its defining formula has no model. The second statement is $\Sat(A) \cap \Sat(B) = \Sat(A \wedge B)$.
\end{proof}

The solver returns either $\mathsf{sat}$, $\mathsf{unsat}$ or $\mathsf{unknown}$. An answer of $\mathsf{sat}$ includes a model, which provides a concrete witness state; $\mathsf{unsat}$ proves emptiness; $\mathsf{unknown}$ establishes nothing and is reported as such.

\begin{example}[Aliasing]
Let $P$ first store $7$ at $p$, then $9$ at $q$, and finally load the byte at $p$. When memory is represented by a single array $M$, the result is
\[
  \mathrm{read}(\mathrm{store}(\mathrm{store}(M,p,7),q,9),p) \;=\;
  \begin{cases} 9 & p = q, \\ 7 & p \neq q, \end{cases}
\]
by the read-over-write axioms. Thus, for any allocation box $B$, the boxes $B \wedge (p=q)$ and $B \wedge (f_P = 9)$ have identical extensions. Proposition~\ref{prop:containment} establishes both inclusions using two unsatisfiability checks.
\end{example}

\begin{proposition}[Projection of irrelevant inputs]
Let $r \in \mathcal{R}$. Then $f_P$ is independent of the cell $r$ if and only if the formula
\[
  f_P(\sigma) \neq f_P(\sigma[r := v])
\]
in the free variables $\sigma, v$ is unsatisfiable.
\end{proposition}

\begin{proof}
Independence of $r$ means $f_P(\sigma) = f_P(\sigma')$ whenever $\sigma$ and $\sigma'$ agree off $r$. Every pair of this kind has the form $(\sigma, \sigma[r:=v])$. The formula is satisfiable precisely when such a pair provides a counterexample.
\end{proof}

Only once this check succeeds is $r$ removed from the input interface.

\section{Execution fibers}

\begin{definition}
Given $f$ from $X$ into $Y$ and $y \in Y$, the \emph{fiber} over $y$ is $f^{-1}(y) = \{x \in X : f(x) = y\}$.
\end{definition}

\begin{example}\label{ex:fiber}
Let $f(x) = 6x + 4$ on $\Z_{256}$. Then $f^{-1}(22) = \{3, 131\}$.
\end{example}

\begin{proof}
$6x + 4 \equiv 22 \pmod{256}$ is $6x \equiv 18 \pmod{256}$. Since $\gcd(6,256) = 2$ divides $18$, this is equivalent to $3x \equiv 9 \pmod{128}$. Because $3 \cdot 43 = 129 \equiv 1 \pmod{128}$, multiplying by $43$ yields $x \equiv 387 \equiv 3 \pmod{128}$. The residues of $\Z_{256}$ that are congruent to $3$ modulo $128$ are $3$ and $131$.
\end{proof}

In general, $|f^{-1}(y)| \in \{0, \gcd(a, 2^w)\}$ for $f(x) = ax + c$ on $\Z_{2^w}$, so the observation by itself cannot distinguish the inputs when $a$ is even.

\begin{proposition}[Augmented observation]
Let $x \in \{0, \dots, 255\}$ and $\tilde f(x) = (y, k)$, where $6x + 4 = y + 256k$ with $0 \le y < 256$. Then $\tilde f$ is injective, and $x = (y + 256k - 4)/6$.
\end{proposition}

\begin{proof}
The integer $6x + 4$ can be recovered from $(y,k)$ as $y + 256k$, while $x$ can be recovered from $6x + 4$ because multiplication by $6$ is injective on $\Z$.
\end{proof}

For the fiber in Example~\ref{ex:fiber}: $6 \cdot 3 + 4 = 22$, $(22, 0)$ follows, while $6 \cdot 131 + 4 = 790 = 22 + 3 \cdot 256$ yields $(22, 3)$. The quotient $k$ is additional observed state: the bits lost when truncating to one byte.

\begin{remark}[Three coordinate systems]
The same function can be represented in three ways: modular arithmetic on $\Z_{256}$; a Boolean circuit over $\F_2^8$ with explicit carry bits; and an embedding into $\mathbb{R}^8$ whose points are constrained to the hypercube vertices $\{0,1\}^8$. Each model has the fiber $\{3,131\}$ over $22$. If the last model is relaxed to the solid cube, the fiber expands to include non-vertex points. Thus, the discreteness constraint belongs to the exact model and cannot be omitted.
\end{remark}

\section{Compiler fibers}

Let $S$ denote a set of source texts, $\theta$ a compiler profile (compiler, version, flags), and $C_\theta$ the map from $S$ to byte strings obtained by compiling and extracting the relevant function or image.

\begin{definition}
For a target byte string $b$, the \emph{compiler fiber} is $C_\theta^{-1}(b) \subseteq S$. A \emph{partial right inverse} on $D$ is a map $I$ from $D$ into $S$ such that $C_\theta(I(b)) = b$ for every $b \in D$.
\end{definition}

Strict source recovery means producing an element of $C_\theta^{-1}(b)$. We distinguish two notions:
\begin{align*}
  \text{exact:} \quad & C_\theta(s) = b, \\
  \text{semantic:} \quad & f_{C_\theta(s)} = f_b \text{ on the modelled domain}.
\end{align*}
Exact recovery implies semantic recovery. The converse does not hold, because different register allocations or instruction selections can compute the same function.

\subsection{The affine family}

\begin{proposition}\label{prop:affine}
Let $f(x,y) = ax + by + c$ on $\Z_{2^w}$. Then
\[
  c = f(0,0), \qquad a = f(1,0) - c, \qquad b = f(0,1) - c,
\]
and distinct triples $(a,b,c) \in \Z_{2^w}^3$ specify distinct functions.
\end{proposition}

\begin{proof}
Evaluate at the three basis points. The formulas recover $(a,b,c)$ from $f$, so any two triples with the same $f$ must coincide.
\end{proof}

Proposition~\ref{prop:affine} provides the inverse $I$. A candidate counts as exact only when recompiling $I(b)$ reproduces the function bytes; otherwise, it is reported as semantic.

\begin{remark}[Non-uniqueness for polynomials]
Uniqueness does not hold beyond degree one. For $w \ge 1$, the product $x(x-1)$ is always even, and therefore
\[
  2^{w-1} x (x - 1) \equiv 0 \pmod{2^w} \quad \text{for all } x.
\]
Thus, $g$ and $g + 2^{w-1}(x^2 - x)$ define the same function on $\Z_{2^w}$ despite having different coefficients. Polynomial functions on $\Z_{2^w}$ form a proper quotient of the polynomial ring, so a recovered coefficient vector is merely a representative of its class.
\end{remark}

\subsection{Whole-program search}

A source grammar has a finite set of holes $H = \{h_1, \dots, h_n\}$, a finite choice set $K_i$ with costs $\kappa_i$ from $K_i$ into $\N$ for each hole, and a finite set of profiles $\Theta$ with cost $\kappa_0$. An assignment selects one choice for each hole, and substitution is literal and applied at every occurrence. The candidate space is therefore
\[
  \Omega = \Theta \times K_1 \times \dots \times K_n, \qquad
  \kappa(\theta, k_1, \dots, k_n) = \kappa_0(\theta) + \sum_{i=1}^n \kappa_i(k_i).
\]
Each $\omega \in \Omega$ specifies a complete project $s_\omega$, and $\omega$ is accepted when $C_\theta(s_\omega) = b$ for every target.

\begin{proposition}\label{prop:order}
Sort each coordinate by cost, then initialize a priority queue with the tuple of minimum choices. Repeatedly pop a minimum-cost tuple, emit it, and push each tuple formed by advancing one coordinate to its next choice, skipping any tuple already seen. This emits every element of $\Omega$ exactly once, in nondecreasing cost order.
\end{proposition}

\begin{proof}
Advancing a coordinate cannot lower the cost, because the choices in each coordinate are sorted. Each tuple other than the initial one has a predecessor formed by moving back one coordinate, with cost no greater than its own. Induction on the sum of the coordinate indices shows that every tuple is pushed, since its predecessor is popped. The seen set prevents duplicates. Suppose a tuple $u$ is popped after a tuple $v$ with $\kappa(u) < \kappa(v)$. When $v$ was popped, some ancestor $u'$ of $u$ along a chain of predecessors was in the queue, with $\kappa(u') \le \kappa(u) < \kappa(v)$, contradicting the minimality of $v$.
\end{proof}

\begin{corollary}
The first accepted candidate has the minimum declared cost among all accepted candidates in $\Omega$. If the budget expires after $N$ candidates, then every rejected candidate whose cost is below the last emitted cost has been tested. No conclusion can be drawn about $\Omega$ beyond those $N$, nor about any sources outside $\Omega$.
\end{corollary}

The size $|\Omega| = |\Theta| \prod_i |K_i|$ grows exponentially in $n$. Ordering makes the search useful, but does not make it polynomial.

\section{Color refinement}

Let $G = (V, E, \ell)$ be a directed multigraph with vertex labels $\ell$ and edge slots $s(e)$ recording operand position; operands of commutative operations share a slot.

\begin{definition}
Set $\chi_0 = \ell$ and
\[
  \chi_{t+1}(v) = \Big( \chi_t(v),\; \big\{\!\big\{ (s(e), \chi_t(u)) : e = (u,v) \big\}\!\big\},\; \big\{\!\big\{ (s(e), \chi_t(u)) : e = (v,u) \big\}\!\big\} \Big),
\]
where $\{\!\{\cdot\}\!\}$ is a multiset. Let $\mathcal{P}_t$ denote the partition of $V$ into classes of $\chi_t$.
\end{definition}

\begin{lemma}
$\mathcal{P}_{t+1}$ refines $\mathcal{P}_t$. If $\mathcal{P}_{t+1} = \mathcal{P}_t$, then $\mathcal{P}_{t+k} = \mathcal{P}_t$ for all $k$. Stabilization takes at most $|V|$ rounds.
\end{lemma}

\begin{proof}
$\chi_{t+1}(v)$ includes $\chi_t(v)$ as its first component, so identical new colors entail identical old colors. If two rounds produce the same partition, then $\chi_{t+1}$ is a bijective relabelling of $\chi_t$ on the classes, and the next round encounters the same multisets up to that relabelling. Each strict refinement increases the number of classes, which is bounded by $|V|$.
\end{proof}

Colors are compared as exact tuples, preventing hash collisions from merging distinct classes.

\begin{proposition}
If $G \cong G'$ under a label- and slot-preserving isomorphism, the stable color histograms of $G$ and $G'$, refined jointly, are equal. Consequently, different histograms prove non-isomorphism.
\end{proposition}

\begin{proof}
For an isomorphism $\varphi$, $\chi_t(\varphi(v)) = \chi_t(v)$ holds for all $t$ by induction on $t$, since it maps the in- and out-neighbourhoods bijectively while preserving slots.
\end{proof}

\begin{example}[The converse fails]
Let $C_6$ be the six-cycle and $2K_3$ be two disjoint triangles, both with uniform labels. Each graph is $2$-regular and has six vertices, so every vertex receives the same color in every round and the histograms match. Nevertheless, the graphs are not isomorphic.

Betti numbers of the underlying $1$-complex distinguish them: $\beta_0$ counts the components, and $\beta_1 = |E| - |V| + \beta_0$ is the cycle rank:
\[
  (\beta_0, \beta_1)(C_6) = (1, 1), \qquad (\beta_0, \beta_1)(2K_3) = (2, 2).
\]
\end{example}

\begin{remark}
Structural matching and semantic equality are independent properties. On $\Z_{2^w}$, $x \cdot 3$ and $(x \ll 1) + x$ represent equal functions despite having different graphs. Conversely, $x + y$ and $x \oplus y$ share the same untyped topology but differ at $x = y = 1$.
\end{remark}

\section{Counterexample-guided synthesis}

Given a finite parameter domain $\mathcal{P}$, a finite state domain $\Sigma_0$, and a violation formula $V(p, \sigma)$, the goal is
\[
  \exists p \in \mathcal{P} \;\; \forall \sigma \in \Sigma_0 : \neg V(p, \sigma).
\]

The procedure maintains a finite set $T \subseteq \Sigma_0$ of counterexamples. It searches for $p$ satisfying $\bigwedge_{\tau \in T} \neg V(p, \tau)$; if none exists, the grammar is exhausted. Otherwise, it searches for $\sigma$ satisfying $V(p, \sigma)$. If none exists, it accepts $p$; if $\sigma$ exists, it adds $\sigma$ to $T$.

\begin{proposition}
Accepted parameters are correct, and the procedure halts within at most $|\mathcal{P}|$ iterations.
\end{proposition}

\begin{proof}
Acceptance is equivalent to the unsatisfiability of $V(p, \sigma)$ over the whole of $\Sigma_0$. If $\sigma$ is added for $p$, then $V(p, \sigma)$ holds, which means $p$ violates the new constraint and cannot be proposed again. Each iteration eliminates at least one parameter.
\end{proof}

Exhaustion establishes only that no parameter in this grammar works. Solver $\mathsf{unknown}$ is reported as unresolved.

\section{Loop contracts}

Let a loop have a one-lap transition $T(\sigma, \sigma')$, extracted from the machine code, a guard $g$, an initial condition $\mathrm{Init}$, an invariant candidate $J$, a postcondition $Q$, and a ranking function $\rho$ into $\N$.

\begin{theorem}
Suppose the following formulas hold:
\begin{align*}
  &\mathrm{Init}(\sigma) \Rightarrow J(\sigma), \\
  &J(\sigma) \wedge g(\sigma) \wedge T(\sigma,\sigma') \Rightarrow J(\sigma'), \\
  &J(\sigma) \wedge \neg g(\sigma) \Rightarrow Q(\sigma), \\
  &J(\sigma) \wedge g(\sigma) \wedge T(\sigma,\sigma') \Rightarrow \rho(\sigma') < \rho(\sigma).
\end{align*}
Then every execution starting from $\mathrm{Init}$ terminates in a state satisfying $Q$.
\end{theorem}

\begin{proof}
Induction on the number of laps shows that $J$ holds at every loop head. During any run, $\rho$ decreases strictly in $\N$ at each lap, so the run has at most $\rho(\sigma_0)$ laps. At exit, $\neg g$ holds along with $J$, and therefore $Q$.
\end{proof}

Each implication is verified by checking that its negation is unsatisfiable. For a counter $i$ that runs to a bound $n$, $\rho = n - i$. A constant $\rho$ fails the fourth condition, producing a counterexample.

\begin{remark}
An XOR accumulator $a' = a \oplus k$ is not affine over $\Z_{2^w}$, but it is linear over $\F_2^w$. Thus, the invariant language depends on the chosen coordinates, as noted in the remark on three coordinate systems.
\end{remark}

\section{Prototype refinement as a fixed point}

Let $F$ denote the set of recovered functions. For each $g \in F$, let $\delta(g)$ be the prototype inferred by the decompiler from the body of $g$: the model, arity, and a storage triple $(\text{space}, \text{offset}, \text{size})$ with a type for each parameter and the return. For a direct call site $c$ in a caller $h$ with target $g$, let $\gamma(c)$ denote the prototype inferred from the $h$ side of the call.

\begin{definition}
The evidence is \emph{consistent} at $c$ if $\gamma(c)$ and $\delta(g)$ match in model, variadicity, custom storage, arity, every storage triple, and every type. A disagreement is a \emph{conflict}.
\end{definition}

Conflicts indicate disagreements between independent inferences; they do not prove anything about the true ABI.

A call is \emph{trusted} if its caller has no unresolved relocation reads, bad data, or unimplemented instructions. For a caller $h$, define the manifest
\[
  m(h) = \{ (g, \delta(g)) : h \text{ calls } g \text{ directly},\; g \neq h,\; \delta(g) \text{ supported} \},
\]
where supported means non-variadic, using standard register or stack storage and supported types. Let $\Phi$ denote the operator that re-decompiles every caller $h$ using the hypotheses $m(h)$ and recomputes $\delta$ and $\gamma$.

The decompiler checks each hypothesis against the storage it reconstructs. If it rejects any hypothesis, or the result contains unresolved machine operations, the previous result for $h$ is retained.

\begin{proposition}
Let $\eta(h)$ be a hash of $m(h)$ combined with the backend identity, and record the last adopted value $\eta^\ast(h)$. If iteration stops because $\eta^\ast(h) = \eta(h)$ holds for every caller $h$, the current state is a fixed point of $\Phi$ in this sense: rerunning $\Phi$ would provide each caller with exactly the hypotheses used in its last decompilation.
\end{proposition}

\begin{proof}
For a deterministic backend, decompiling $h$ depends only on $h$, the backend, and $m(h)$. Assuming there are no hash collisions, equal hashes imply equal manifests, so $\Phi$ reproduces the current evidence.
\end{proof}

No lattice order on prototypes is claimed under which $\Phi$ is monotone, so convergence is not generally guaranteed. Termination is enforced by a round limit $R$, a time budget, and a rule permitting each pair $(h, \eta)$ to be attempted at most once. The outcome is reported as one of: fixed point, unresolved (a round with no adoption), or budget exhausted.

\begin{example}
Let $\mathtt{leaf}(x,y) = x + y$ take $x$ in $\mathrm{EDI}$ and $y$ in $\mathrm{ESI}$, and suppose the caller sets only $\mathrm{EDI} = 7$ before the call. In its first decompilation, the caller sees one argument. With $m(\mathtt{caller}) = \{(\mathtt{leaf}, \delta(\mathtt{leaf}))\}$, the call becomes binary, and the forwarded $\mathrm{ESI}$ appears as a caller parameter. A hypothesis assigning the first argument to the wrong register is rejected because of the storage mismatch. On rerun, $\eta^\ast = \eta$ and no reanalysis occurs; corrupting $\eta^\ast$ triggers exactly one reanalysis, which restores the same $\eta$.
\end{example}

\section{Translation unit fibers}

At high optimisation levels, inlining, cloning and interprocedural register allocation couple the functions in a translation unit. This section examines how much the compiler fiber still decomposes into per-function pieces.

\subsection{Summary locality}

Express the source space as a product $S = S_1 \times \dots \times S_m$ with one slot for each function, grouping functions that are inlined into one another into a shared slot. Let $N(j)$ denote the call-graph neighbours of slot $j$. Each slot has a finite summary set $A_j$, which records facts such as purity, the noreturn property, clobbered registers, propagated constants and memory effects. Set $A = A_1 \times \dots \times A_m$ and denote by $a_{N(j)}$ the restriction of $a \in A$ to the neighbours of $j$.

\begin{definition}[Summary locality]
There are maps $\sigma_j$ from $S_j \times \prod_{i \in N(j)} A_i$ to $A_j$ and $F_j$ from $S_j \times \prod_{i \in N(j)} A_i$ to byte strings such that, for every $s \in S$, the system
\[
  a_j = \sigma_j(s_j, a_{N(j)}), \qquad j = 1, \dots, m,
\]
has exactly one solution $a(s) \in A$, and the bytes of function $j$ are $C(s)_j = F_j(s_j, a(s)_{N(j)})$.
\end{definition}

\begin{definition}
For a target $b = (b_1, \dots, b_m)$ and $a \in A$, define the \emph{local fiber} of slot $j$ by
\[
  L_j(a) = \{ s_j \in S_j : \sigma_j(s_j, a_{N(j)}) = a_j \text{ and } F_j(s_j, a_{N(j)}) = b_j \}.
\]
\end{definition}

\begin{theorem}\label{thm:tu-product}
Under summary locality,
\[
  C^{-1}(b) = \bigsqcup_{a \in A} \; \prod_{j=1}^m L_j(a).
\]
\end{theorem}

\begin{proof}
Let $s \in C^{-1}(b)$ and set $a = a(s)$. Then $a$ solves the system and $F_j(s_j, a_{N(j)}) = b_j$ for every $j$, hence $s_j \in L_j(a)$. Conversely, suppose $s_j \in L_j(a)$ for every $j$. The first condition in $L_j(a)$ implies that $a$ solves the system for $s$, so uniqueness gives $a = a(s)$, and the second condition then yields $C(s)_j = b_j$. The union is disjoint since $a(s)$ is a function of $s$.
\end{proof}

\begin{corollary}
For a fixed $a$, searching the product requires $\sum_j |\Omega_j|$ candidate checks rather than $\prod_j |\Omega_j|$, where $\Omega_j$ is the candidate set for slot $j$. A summary is \emph{observable} if it is determined by $b$; examples include a clobber set read from the callee bytes or a constant-propagation clone exposed as a symbol. If all summaries are observable, only one $a$ is possible, and the search is linear in $m$.
\end{corollary}

\begin{proposition}
Suppose some summaries are unobservable, the call graph has treewidth $w$, and $|A_j| \le \alpha$ for every $j$. Then deciding whether $C^{-1}(b)$ is nonempty and finding one of its elements take $O(m \cdot \alpha^{w+1} \cdot \max_j |\Omega_j|)$ local checks.
\end{proposition}

\begin{proof}
Every constraint in the system and every byte condition involves only slot $j$ and its neighbours. Thus, the constraints form a network over variables $a_j$ whose primal graph is the call graph. Take a tree decomposition of width $w$. Elimination along this decomposition enumerates at most $\alpha^{w+1}$ assignments for each bag, and tests the slots in each assignment against at most $\max_j |\Omega_j|$ candidates. The decomposition has $O(m)$ bags.
\end{proof}

For fixed $w$ and $\alpha$, the bound is polynomial in the number of functions. It offers no improvement for $|\Omega_j|$, which remains exponential in the number of holes in a single function.

\begin{remark}[Refutation test]
Summary locality predicts that changing slot $j$ while keeping $a$ fixed leaves $C(s)_i$ unchanged for $i \neq j$. A change to function $i$ can therefore be explained only if it occupies slot $j$, neighbours a changed function in the call graph, or has a changed inlining decision.
\end{remark}

\subsection{Global inlining}

gcc's interprocedural inliner greedily processes call edges from a priority queue, $e$, ordered by badness $\beta_e$. Each edge is also subject to per-caller limits, represented by an indicator $\lambda_e \in \{0, 1\}$, and a unit budget
\[
  B(s) = (1 + p) \max(N(s), L), \qquad N(s) = \sum_j n_j(s_j),
\]
where $p$ is the inline-unit-growth parameter, $L$ is large-unit-insns, and $n_j$ is the size estimate for slot $j$. For a fixed processing order, an edge $e$ with growth $g_e$ is included in the decision set $D(s)$ exactly when
\[
  \lambda_e = 1 \quad \text{and} \quad N(s) + \sum_{f \in D(s),\, f \prec e} g_f + g_e \le B(s).
\]

\begin{proposition}
Extend summary locality so that $\sigma_j$ and $F_j$ may also depend on $D$, and assume that $D(s)$ is a function of $s$. Then
\[
  C^{-1}(b) = \bigsqcup_{(a, D)} \Big\{ s \in \prod_{j=1}^m L_j(a, D) : D(s) = D \Big\}.
\]
\end{proposition}

\begin{proof}
Theorem~\ref{thm:tu-product} applies by the same argument, with $(a, D)$ replacing $a$. The only condition that does not decompose over slots is $D(s) = D$, which remains as a filter.
\end{proof}

For a fixed order, the filter is a set of linear inequalities over the integers $n_j$ and $g_e$. It is conjectured, but not proved, that a dynamic program over the slots that tracks prefix sums can decide the filter in pseudo-polynomial time in $B$. Two issues remain open. The order compares $\beta_e$ across all edges, allowing up to $|E|!$ possible orders, while gcc recalculates $g_e$ after each inlining and $B$ uses the original size. It is unknown whether the entire order must be fixed, or whether it suffices to fix, for each rejected edge, the set of edges preceding it.

\subsection{Indicator functions and spectra}

Represent a set of sources by its indicator function. Over $\F_2$, Theorem~\ref{thm:tu-product} becomes
\[
  \mathbf{1}_{C^{-1}(b)} = \sum_{a \in A} \prod_{j=1}^m \mathbf{1}_{L_j(a)},
\]
This holds at every $s$ where the system has a unique solution, and may fail otherwise. The $\F_2$ affine hull of a set gives parity invariants through Gaussian elimination, but is not the same as the real affine hull. For example, $X = \{000, 110, 101, 011\}$ is an $\F_2$ hyperplane but has full affine rank over $\mathbb{R}$.

Assign each slot a graph $G_j$ on $S_j$, and take their Cartesian product $G = G_1 \,\square\, \dots \,\square\, G_m$. Its Laplacian eigenvectors are tensor products of eigenvectors from the factors, with eigenvalues given by addition. For the hypercube $Q_n = K_2 \,\square\, \dots \,\square\, K_2$, the Walsh characters $\chi_T$ have eigenvalue $2|T|$. Over $\mathbb{R}$, the spectral coefficient of $\mathbf{1}_{C^{-1}(b)}$ at $T = (T_1, \dots, T_m)$ is $\sum_a \prod_j$ of the coefficient of $\mathbf{1}_{L_j(a)}$ at $T_j$, again whenever $a(s)$ is unique.

\begin{proposition}
Let $\mathcal{L}$ be the Laplacian of a finite graph, with orthonormal eigenvectors $v_k$ and eigenvalues $\lambda_k$. Let $X$ be a vertex set and write $\mathbf{1}_X = \sum_k c_k v_k$. Then
\[
  \sum_k \lambda_k c_k^2 = |\partial X|,
\]
is the number of edges with exactly one endpoint in $X$. On $Q_n$, taking $\hat f(T) = 2^{-n} \sum_x f(x) (-1)^{T \cdot x}$, this becomes $|\partial X| = 2^{n+1} \sum_T |T| \, \hat{\mathbf{1}}_X(T)^2$.
\end{proposition}

\begin{proof}
$\mathbf{1}_X^\top \mathcal{L} \mathbf{1}_X = \sum_{\{x,y\} \in E} (\mathbf{1}_X(x) - \mathbf{1}_X(y))^2$, where each summand is $1$ for a cut edge and $0$ otherwise. On $Q_n$, the characters are orthogonal with $\|\chi_T\|^2 = 2^n$, so $c_T^2 = 2^n \hat{\mathbf{1}}_X(T)^2$ and $\lambda_T = 2|T|$.
\end{proof}

\subsection{Slot by slot search}

First, read the relocations and call targets. These determine the call graph and all observable summaries before the search begins. Then restrict the slots one by one to their local fibers, following an elimination order along the call graph. Each step checks an exact equality, so no point in the fiber is discarded; the cost depends on the width of the order. This is sound only if the summaries fixed in advance are genuinely functions of $b$. Uniqueness of $a(s)$ is the sole hypothesis that neither reading the references first nor ordering the slots can provide.

\subsection{Locality on real gcc}

The refutation test used gcc 15.2 at -O3 on two zstd 1.5.6 translation units. Compiling the preprocessed units reproduces the shipped object code for 20 of 20 functions in huf\_compress and 147 of 147 in zstd\_compress. In turn, a probe—a single store to a static volatile variable—was inserted at the start of each definition, and every function was compared with relocations masked.

\begin{itemize}
  \item huf\_compress: 125 changes in the probed slot, 11 at neighbours, 0 from inlining and 4 unexplained. The 4 lie two calls away from the probe, across a function whose bytes do not change, and disappear under -fno-ipa-modref. A memory effect summary therefore passes through a function without changing its bytes. This is consistent with summary locality, but the summary is not observable.
  \item zstd\_compress: 1146 changes in the probed slot, 174 at neighbours, 296 from inlining, 22 unexplained and 31 added or removed symbols, mostly part and cold splits. Of the inlining changes, 288 come from one call site in ZSTD\_CCtxParams\_setParameter that sits on an inlining threshold and flips with almost any size change elsewhere.
\end{itemize}

Global coupling does occur, but it is concentrated in a few marginal call sites, leaving $D$ with few free bits in practice. Because a probe deliberately introduces a side effect, these counts may exaggerate the impact of an edit that preserves the summaries.

\section{Thermodynamics of destroyed information}

Compilation discards information. The sources
\[
  \texttt{int f(void) \char123{} return 42; \char125} \qquad \text{and} \qquad \texttt{int f(void) \char123{} int x = 999; return 42; \char125}
\]
compile to identical bytes and therefore belong to the same fiber; nothing computed from those bytes can identify which source was written. This section models the fiber as the ground state of a statistical system and asks which aspects of the source remain.

\subsection{Gibbs measures on the source space}

Fix a target $b$, a distance $d$ between byte strings that is $d(x, y) = 0$ only for $x = y$, and a graph $G$ on $S$ with edges representing single edits. For $s \in S$, define
\[
  u(s) = \#\{ s' : s' \sim s \text{ in } G,\; C(s') \neq C(s) \},
\]
as the number of neighbours that compile differently, and let $|s|$ denote the size of $s$. The energy is
\[
  E(s) = d(C(s), b) + \mu\, u(s) + \lambda |s|,
\]
and the Gibbs measure at temperature $T > 0$ is $p_T(s) = Z_T^{-1} \exp(-E(s)/T)$.

The value $u(s)$ depends only on the neighbourhood of $s$, and
\[
  \sum_{s \in X} u(s) = |\partial X| \quad \text{for } X = C^{-1}(C(s)),
\]
Thus, $u$ is the pointwise version of the boundary count.

\begin{remark}[Where the topology enters]
$p_T$ depends on $G$ only through $u$. Without the term $\mu u$, the equilibrium measure is identical for every graph on $S$; the graph then influences only the dynamics. A Metropolis chain on a connected $G$ with symmetric proposals has $p_T$ as its stationary law, while $G$ determines its mixing rate and where it becomes trapped. In what follows, the dynamics encodes the global topology and the energy encodes the local topology.
\end{remark}

\begin{remark}[Size]
$S$ is infinite because the dead store above can be repeated with any constant. Consequently, $Z_T$ diverges unless there is a size term or a size cap. With $\lambda > 0$, the limit $T \to 0$ selects the shortest elements of the fiber, treating destroyed parts as absent, as in a minimum description length principle.
\end{remark}

\subsection{Soundness of the low temperature limit}

\begin{proposition}\label{prop:scaling}
Let $S$ be finite, $C^{-1}(b) \neq \emptyset$, and define $\mu = \rho T$ and $\lambda = \ell T$ using constants $\rho, \ell \ge 0$. Then, as $T \to 0$, the measure $p_T$ converges to
\[
  p_0(s) \propto \mathbf{1}[C(s) = b] \, \exp(-\rho\, u(s) - \ell |s|),
\]
whose support is the entire fiber. If instead $\mu > 0$ is fixed, $p_T$ concentrates on the minimisers of $d(C(s), b) + \mu u(s)$. In general, these do not include every point in the fiber and may exclude all of them.
\end{proposition}

\begin{proof}
Under this scaling, $\exp(-E(s)/T) = \exp(-d(C(s), b)/T) \exp(-\rho u(s) - \ell |s|)$. The second factor is independent of $T$ and is positive. The first is $1$ on the fiber and no greater than $\exp(-\delta/T)$ elsewhere, where $\delta > 0$ is the smallest nonzero value of $d$ on the finite set $C(S)$. If $\mu$ is fixed, the exponent $-E/T$ reaches its maximum precisely at the minimisers of $E$, and these points carry all the mass in the limit.
\end{proof}

\subsection{What can be recovered}

Write $s \approx s'$ if $C(s) = C(s')$. The destroyed part of $s$ is its equivalence class in $S/{\approx}$, and its size up to $n$ is quantified by
\[
  H_n(b) = \log_2 \big| C^{-1}(b) \cap \{ s : |s| \le n \} \big|.
\]
The data processing inequality says that no function of $b$ contains more information about $s$ than $b$; therefore, the conditional entropy $H(s \mid b)$ is a hard lower bound. The terms $\mu u$ and $\lambda |s|$ select a representative of the class. They serve as a prior and recover no information from the binary.

Erasure at -O3 is not invariably complete. Dead code that remains until the interprocedural passes is still included in size estimates and summaries, and can shift a marginal inlining decision in another function.

\begin{definition}
Consider $s \approx s'$ at the single-function level. The pair is \emph{tier 1} if $s$ and $s'$ remain equivalent in every translation-unit context; it is \emph{tier 2} if some context distinguishes them. Sources that compile differently on their own form a tier 3 pair.
\end{definition}

The experiments below disprove the conjecture that tier 2 information affects the bytes only through $D$, and point to the following hypothesis.

\begin{remark}[Conjecture]
In a translation-unit context, $H(s \mid b_{\mathrm{TU}}) = H(s \mid b_f, a, D)$, where $b_f$ denotes the bytes of the function itself and $(a, D)$ its extended summary.
\end{remark}

\subsection{Global topology of the fiber}

The fiber may divide into several components within $G$. At low temperature, a chain can move between components only by crossing an energy barrier $\Delta$, which takes time of order $\exp(\Delta/T)$. The relevant invariant is the component structure of the sublevel sets $\{ s : E(s) \le t \}$ as $t$ changes—in other words, the persistence of $H_0$. A fiber component born at level $0$ and merging at level $\Delta$ forms a finite bar $[0, \Delta)$. The chain's spectral gap and the Cheeger constant of $G$ weighted by $p_T$ bound the mixing time.

\subsection{Experiments}

All experiments were run with gcc 15.2 at -O3. They are small, exhaustive tests, not measurements on zstd.

\paragraph{Dead code fibers.} Four base functions—a constant return, a loop sum, a call wrapper and a branch—were given up to two insertions selected from fifteen kinds of dead code, yielding 964 sources in total. Every source compiles to the bytes of its base function, so each base fiber contains 241 points and $H_2 = 7.91$ bits. Twelve of the fifteen kinds, among them dead loops, dead arrays and branches rendered dead by inlining a helper, leave no trace even in the interprocedural size summary and are therefore tier 1. The exceptions are branches guarded by a static variable that is never written. They retain their size in the summary because the compiler only discovers that the variable is never written during the interprocedural passes. They account for $1.69$ of the $7.91$ bits, and a caller close to the inlining threshold recovers between $1.62$ and $1.68$ of those bits. None are recovered for the constant base, which is always inlined.

\paragraph{Exact thermodynamics.} The loop-sum base, with the fifteen dead kinds and five live edits and at most two insertions, produces $421$ sources in $27$ byte classes. Inserting, deleting or replacing one insertion defines $G$, whose maximum degree is $59$; $d$ is the byte edit distance.
\begin{itemize}
  \item Proposition~\ref{prop:scaling} is borne out. For the base target, whose fiber has $242$ points, scaling $\mu = \rho T$ retains all $242$ in the limit, whereas every tested fixed $\mu$ retains exactly one. For a second target, the fixed value $\mu = 10$ shifts the ground state entirely outside the fiber.
  \item The translation unit reveals more than $D$. With a caller at three inlining thresholds, $H(s \mid b_{\mathrm{TU}})$ is lower than $H(s \mid b_f, D)$ by between $0.26$ and $0.83$ bits. In one case, the function bytes are identical and both calls are inlined, yet the caller is $321$ bytes long in one instance and $501$ in the other. The dead guard makes the callee impure from the caller's perspective, preventing the caller from merging its two calls. This information is in a summary in $a$, not in a bit of $D$.
  \item The base fiber has two components in $G$: a large component with $241$ points and an isolated point, giving one finite bar $[0, 33)$. At the isolated point, the live edit \texttt{a \^{}= 3} is applied twice. The edits cancel, amounting to algebraic erasure rather than dead code, and every neighbour of that point compiles differently. The chain's spectral gap drops from $7.6 \times 10^{-2}$ at $T = 30$ to $5.4 \times 10^{-7}$ at $T = 3$, and is numerically zero below $T = 1$. For the second target, the fiber is connected, has no finite bar, and the gap remains at $8.5 \times 10^{-3}$ as $T \to 0$.
\end{itemize}

\paragraph{Recovery.} Taking each of the $421$ sources in turn as the hidden truth, the most probable element under the limit measure always compiles to the target bytes. It matches the hidden source in $6.4$ percent of cases when only the function bytes are observed, and in $9.5$ percent when the translation unit is observed. It preserves the same live edits in $91.4$ and $92.9$ percent of cases. The mean remaining entropy is $6.33$ and $5.22$ bits, respectively. An annealed Metropolis chain, run $1200$ times from random starting points for $4000$ steps each, ends at a source with the target bytes in $99.5$ percent of runs.

These figures describe a toy space for one function and cannot, on their own, predict recovery from a real binary, where finding any candidate in the first place is the hard part.
\section{Limitations}

All results in these notes are proved with respect to the machine model $\Sigma$ and a supplied, rather than inferred, ABI. Accordingly, box containment, the fibers of Section 3 and the loop invariants of Section 7 make no claims about faults, page permissions, concurrency or timing. They hold only if the binary actually uses the stated calling convention.

Recovery is judged by the strict standard $C_\theta(s) = b$. A candidate that computes the same function over the modelled domain is classified as semantic agreement, never as recovered source, since different register allocations and instruction selections routinely implement the same function. The reverse gap is equally important: if the search over $\Omega$ runs out of budget, the sole conclusion is that the candidates tested failed. This does not show that no source exists, either within the grammar or outside it.

The structural tools have the same one-sided character. Different stable 1-WL histograms prove that two graphs are not isomorphic, whereas matching histograms prove nothing, as $C_6$ and $2K_3$ show; moreover, graph isomorphism neither implies nor rules out semantic equality. Finally, a prototype fixed point certifies only that inference has stopped changing. Since no monotone order is known for $\Phi$, convergence is enforced through budgets rather than guaranteed, and a stable set of prototypes does not prove that the recovered program is equivalent to the original.

\end{document}
